% TOP_PROBLEM: {"id": 3218, "title": "Bouchet's 6-flow conjecture", "category_id": 3, "category": {"id": 3, "name": "graph_theory", "display_name": "Graph Theory", "description": "Problems involving graphs, networks, and their properties.", "slug": "graph-theory", "order_index": 3, "created_at": "2026-07-31T15:26:25.671Z"}, "status": "open", "problem_number": "OPG-131", "set_id": 12}
% FIRST_POSTED: 2026-10-01
% ATTEMPT_STATUS: unresolved
% MODEL: GPT 6 Astra Ultra
% UnsolvedMath ID 3218; source catalogue code OPG-131
% !TeX program = lualatex
% Research attempt dated 2026-10-01; canonical statement preserved.
\documentclass[11pt,a4paper]{article}
\usepackage[margin=25mm]{geometry}
\usepackage{fontspec}
\setmainfont{lmroman10-regular.otf}[BoldFont=lmroman10-bold.otf,
 ItalicFont=lmroman10-italic.otf,BoldItalicFont=lmroman10-bolditalic.otf]
\usepackage{amsmath,amssymb,mathtools}
\usepackage{unicode-math}
\setmathfont{latinmodern-math.otf}
\AtBeginDocument{\let\setminus\smallsetminus}
\usepackage{xurl,hyperref}
\hypersetup{colorlinks=true,urlcolor=blue}
\setcounter{secnumdepth}{0}
\setlength{\parindent}{0pt}
\setlength{\parskip}{5pt}
\setlength{\emergencystretch}{3em}
\begin{document}
\section*{Bouchet's 6-flow conjecture}\label{3218}
{\small UnsolvedMath ID 3218 · OPG-131 · Graph Theory · OpenGarden}
\subsection{Definitions and mathematical statement}
A finite bidirected multigraph consists
of an undirected multigraph together
with a sign on each incidence,
or half-edge, at a vertex. Use
sign $+1$ for an end pointing
into its vertex and $-1$ for
an end pointing out. The two
ends of an edge are chosen
independently: both may point in,
both out, or in opposite directions.
A loop has two incidences at
the same vertex, each counted
separately in the sums below.

For an integer $q\geq2$, an integer
$q$-flow is a map $\phi:E\to\mathbb Z$
with $|\phi(e)|<q$ and
\[
 \sum_{h\text{ incident with }v}
       \sigma(h)\phi(e(h))=0
 \qquad\text{for every vertex }v,
\]
where $e(h)$ is the edge of
the half-edge $h$. The flow
is nowhere-zero if $\phi(e)\neq0$
for every edge. A bidirected
graph is flow-admissible if it
has a nowhere-zero integer
$q$-flow for some integer $q$.

Bouchet's conjecture states that
every flow-admissible finite
bidirected graph has a nowhere-zero
integer $6$-flow, with values
in $\{-5,-4,-3,-2,-1,1,2,3,4,5\}$.
The bidirected incidence signs
are fixed as part of the
input and must be respected
by the smaller flow.

The premise is existence of
an integer flow of some size,
not just absence of bridges
in the underlying undirected
graph. Unlike ordinary orientations,
an edge can contribute with
the same sign at both endpoints;
the half-edge definition includes
this distinction explicitly.
\subsection{Short English statement}
Whenever a bidirected graph admits any nonzero integer flow, can it admit one with every edge magnitude at most five?
\subsection{Research attempt}
\paragraph{Scope and method.}
This research attempt by GPT-6 Astra Ultra is dated 1 October 2026.
It preserves the catalogue statement and distinguishes elementary proofs
below from cited prior work. No novelty claim or formal proof-assistant
verification is made. The final paragraph states the precise remaining scope.
\paragraph{Literature and the incidence convention.}
DeVos [R1] proves a universal twelve-flow bound for the same bidirected
integer-flow framework. The present attempt checks elementary cases and
obstructions at the sharper six-flow target. Write $B$ for the incidence
matrix whose column contains the two half-edge signs; for a loop they
are added in the same row. Conservation is exactly $B\phi=0$ over
$\mathbb Z$. A loop with equal half-edge signs contributes $\pm2\phi(e)$,
so treating every loop as a zero column would change the problem.

\paragraph{Switching and balanced inputs.}
Multiplying one row of $B$ by $-1$ reverses all half-edges at that vertex
and preserves its integer kernel. Multiplying an edge column by $-1$
corresponds to replacing the flow value on that edge by its negative,
preserving magnitudes. Give a non-loop edge sign
$s_e=-\sigma(h_1)\sigma(h_2)$, so positive edges have opposite incidences.
Vertex switching changes the sign of every non-loop edge incident there.

Suppose the product of edge signs is positive around every cycle, with
loops interpreted as one-edge cycles. Choose a spanning tree in each
component and switch vertices recursively to make its edges positive.
Every remaining edge then is positive as well: its fundamental cycle
has positive product and all the other edges on it are positive.
Positive loops already have opposite incidences. Thus a balanced input
can be switched to an ordinary directed graph. If its underlying graph
has only even degrees, orient Euler tours and assign unit values.
Edge-column reversals convert this convenient orientation to the switched
input orientation; inverse row switching returns to the original input.
This proves a nowhere-zero two-flow for balanced Eulerian inputs, while
respecting the original fixed half-edge signs through the transformations.

\paragraph{A bridge can carry a nonzero signed flow.}
Take vertices $u,v$ connected by a path of any positive length. Put a
loop with two incoming half-edges at $u$ and a loop with two outgoing
half-edges at $v$. Direct the path from $u$ to $v$. Assign value one
to each loop and value two to every path edge. At $u$ the equation is
$2\cdot1-2=0$; at $v$ it is $2-2\cdot1=0$; all internal path vertices
have one incoming and one outgoing value two. This is a nowhere-zero
three-flow even though every path edge is an underlying bridge. The
example illustrates precisely why ordinary bridge arguments cannot be
imported into the signed conjecture.

\paragraph{An unbalanced cycle obstruction.}
A single cycle with negative sign product has no nowhere-zero integer
flow. At each degree-two vertex conservation makes the two incident
values equal up to sign. Propagating these equations around the cycle
returns $\phi(e)=-\phi(e)$, so $\phi(e)=0$ over the integers and then
every edge value is zero. For a single negative loop this is just
$2\phi(e)=0$. The corresponding modular equation can have nonzero
solutions in a group with two-torsion, demonstrating that group flows
cannot be substituted for the requested integer flow without a lifting
argument. These obstructed graphs fail the flow-admissibility premise;
they are not counterexamples to Bouchet's conjecture.

\paragraph{Where summing signed circuits fails.}
The path-and-loops example contains unavoidable magnitude two on its
connecting path. Summing many such elementary circulations can cancel
on common edges or accumulate values exceeding five. Choosing rapidly
increasing weights prevents some cancellations but gives no fixed upper
bound. A proof for arbitrary admissible inputs must coordinate overlaps
and signs; the balanced Euler-tour proof removes exactly the obstruction
that makes the general case difficult.

\paragraph{Verification and final status.}
The root batch script builds the signed incidence matrices for path
lengths one through twenty and verifies the explicit one/two certificates
over the integers. It checks both loop incidences separately.
\textbf{UNRESOLVED.} Balanced Eulerian inputs and the signed handcuff
family above satisfy the target. We have not reduced an arbitrary
flow-admissible signed graph to these cases while keeping all magnitudes
at most five.

\subsection{Sources}
\begin{itemize}
\item[S1] ULAM.AI and the UnsolvedMath Contributors, supplied problems.json, record 3218 (OPG-131), OpenGarden collection. Catalogue identity and source transcription; CC BY 4.0.
\url{https://huggingface.co/datasets/ulamai/UnsolvedMath}.
\item[S2] Open Problem Garden, Bouchet's 6-flow conjecture. Original problem and discussion; the mathematical formulation below is expanded from the supplied source record.
\url{http://www.openproblemgarden.org/op/bouchets_6_flow_conjecture}.
\item[R1] Matt DeVos, \emph{Flows on Bidirected Graphs},
arXiv:1310.8406. Primary abstract checked 1 October 2026.
\url{https://arxiv.org/abs/1310.8406}.
\end{itemize}
\end{document}
